% Source: https://artofproblemsolving.com/community/c3h2613625_exponents_equal
% Stable identity: aops-topic-2613625-q01
\textbf{(a)} We use the basic fact that if $c>1$ and $t>0$, then $c^t>1$. (For a positive integer $t$ this follows by repeated multiplication. For $t=\dfrac pq$ it follows because $\left(c^{\dfrac pq}\right)^q=c^p>1$, and a positive number at most $1$ has all its powers at most $1$. Real exponents are defined so that this property is kept.)

Suppose $c^x=c^y$ with $x\ne y$, say $t=x-y>0$ (otherwise swap $x$ and $y$). Dividing, $c^t=1$.
\begin{itemize}
\item If $c>1$, this contradicts $c^t>1$.
\item If $0<c<1$, then $\dfrac1c>1$ and $\left(\dfrac1c\right)^t=\dfrac{1}{c^t}=1$, again a contradiction.
\end{itemize}
So $x=y$. $\blacksquare$

\textbf{(b)} For $c=1$, $1^2=1^5=1$ but $2\ne5$. $\blacksquare$

(For $c\le0$ real powers $c^x$ are not defined in general, which is why $c>0$ is assumed.)
